%- block legend
\begin{center}
\includegraphics[scale=0.6]{mario.png}
\end{center}
Odracir é um grande fã de jogos retrô e de computação. Em uma tarde, ele estava jogando seu jogo favorito, \textit{Super Mario Bros.}, em seu Nintendinho.

No jogo, Mario pode estar em dois tamanhos: Alto ou Baixo. Quando Mario está Alto e recebe dano de um inimigo, ele diminui de tamanho e passa a ficar Baixo. Caso Mario já esteja Baixo ao receber dano, ele morre.

Mario também pode encontrar um cogumelo, que faz com que ele cresça. Assim, caso esteja Baixo, ele passa a ficar Alto. Se já estiver Alto, o cogumelo não altera seu tamanho.

Além disso, Mario pode encontrar uma estrela, que o torna invulnerável a danos pelos próximos 5 obstáculos. Enquanto estiver invulnerável, inimigos não causam dano a Mario, mas os obstáculos ainda são contabilizados.

Certo dia, devido a um curto-circuito em seu console, o \textit{save} do jogo de Odracir foi corrompido. Agora ele não sabe qual foi o destino de Mario na fase em que estava. Felizmente, usando suas habilidades computacionais, conseguiu recuperar parte do arquivo corrompido, contendo apenas a sequência de obstáculos que Mario encontraria ao longo da fase.

Cada obstáculo pode ser:

\begin{itemize}
\item um inimigo, que causa dano a Mario;
\item um cogumelo, que faz Mario crescer caso ele esteja Baixo;
\item uma estrela, que torna Mario invulnerável a danos aos próximos 5 obstáculos.
\end{itemize}
Considere que Mario começa a fase Alto e que os obstáculos são encontrados na ordem em que aparecem na sequência.

Dada uma string representando todos os obstáculos da fase, determine se Mario morrerá ou sobreviverá à fase. Caso sobreviva, determine também se ele terminará a fase Alto ou Baixo.
%- endblock

%- block input
A entrada consiste de apenas uma string $S$ representando os obstáculos da fase com pelo menos $1$ e no máximo $10^5$ letras, cada letra pode ser $D$, $C$ ou $E$, representando um obstáculo de \textit{Dano}, \textit{Cogumelo} e \textit{Estrela} respectivamente.
%- endblock

%- block output
A saída deve ser uma linha com uma das seguintes 3 palavras: ``MORTO'' caso Mario morra no percurso, ``ALTO'' caso Mario sobreviva com sua estatura inicial, ou ``BAIXO'' caso Mario sobreviva com a estatura menor que sua inicial.
%- endblock

%- block explanation_0
Observe um fluxograma representando a mudança no estado do Mario a cada obstáculo do exemplo:

\[
\begin{aligned}
\text{Alto} \overset{\texttt{D}}{\rightarrow} \text{Baixo}
\overset{\texttt{C}}{\rightarrow} \text{Alto}
\overset{\texttt{C}}{\rightarrow} \text{Alto}
\overset{\texttt{E}}{\rightarrow} \text{Alto}
\overset{\texttt{D}}{\rightarrow} \text{Alto}
\overset{\texttt{C}}{\rightarrow} \text{Alto}
\overset{\texttt{E}}{\rightarrow} \text{Alto}
\overset{\texttt{D}}{\rightarrow} \text{Alto}
\\
\overset{\texttt{D}}{\rightarrow} \text{Alto}
\overset{\texttt{C}}{\rightarrow} \text{Alto}
\overset{\texttt{D}}{\rightarrow} \text{Alto}
\overset{\texttt{C}}{\rightarrow} \text{Alto}
\overset{\texttt{D}}{\rightarrow} \text{Baixo}
\overset{\texttt{C}}{\rightarrow} \text{Alto}
\end{aligned}
\]
%- endblock

%- block explanation_2
    Neste exemplo no terceiro obstáculo Mario já está baixo e recebe dano, logo ele morre neste momento.
%- endblock

%- block notes
Caso Mario morra ele permanece neste estado mesmo que ainda há mais obstáculos a frente de onde ele recebeu dano.
%- endblock

%- block editorial
Podemos implementar uma solução iterando um loop por todas letras da string em ordem e contando qual a estatura atual do Mario, se no estado atual ele está sobre efeito de uma estrela, e caso sim, contar também contabilizar a duração do efeito, neste caso durante os obstáculos
sobre efeito da estrela ignoramos todos os danos.

Complexidade: $O(N)$ - Onde $N$ é o tamanho da string.
%- endblock
